\bumppoints{20}

\question
The rabbitsfoot mussel was once widespread throughout rivers of eastern North America but now remains in only a few very small populations. The decrease occurred rapidly due to habitat loss. Genetics studies show that the remaining populations have little or no genetic variation. The best explanation for this result is:

\begin{choices}
\choice Inbreeding caused an increase in heterozygosity. 
\choice Balancing selection favored only one or two alleles. 
\correctchoice Genetic drift caused the loss of heterozygosity. 
\choice A bottleneck caused the loss of homozygosity. 
\choice Directional selection favored only one or two alleles. 
\end{choices}

\question
Which \emph{one} of the following statements is true of natural selection?

\begin{choices}
\correctchoice It requires genetic variation, results in descent with modification, and involves differences in reproductive success. 
\choice It requires genetic variation. 
\choice It results in descent with modification. 
\choice It involves differences in reproductive success. 
\choice It results in descent with modification and involves differences in reproductive success. 
\end{choices}

\question
\textit{Anopheles} mosquitoes, which carry the malaria parasite, cannot live above elevations of 5,900 feet. In addition, oxygen availability decreases with higher altitude. Consider a hypothetical human population that is adapted to life on the slopes of Mt. Kilimanjaro in Tanzania, a country in equatorial Africa. Mt. Kilimanjaro's base is about 2,600 feet above sea level and its peak is 19,341 feet above sea level. If the frequency of the sickle-cell allele in the population is plotted against altitude (feet above sea level), which of the following distributions is most likely, assuming little migration of people up or down the mountain?

\begin{multicols}{2}
\begin{choices}
\correctchoice \includegraphics[width=0.35\textwidth]{anopB} 
\choice \includegraphics[width=0.35\textwidth]{anopA} 
\choice \includegraphics[width=0.35\textwidth]{anopC} 
\choice \includegraphics[width=0.35\textwidth]{anopD} 
\end{choices}
\end{multicols}

\newpage

\question
Which mouse has the greatest relative fitness? (A clutch is the offspring from one breeding event.)

\begin{choices}
\choice One that lives for 2 years, has 3 clutches per year, and 8 offspring per clutch. 
\choice One that lives for 4 years, has 3 clutches per year, and 4 offspring per clutch.
\choice One that lives for 3 years, has 3 clutches per year, and 5 offspring per clutch.
\correctchoice One that lives for 3 years, has 3 clutches per year, and 7 offspring per clutch. 
\choice One that lives for 3 years, has 2 clutches per year, and 6 offspring per clutch. 
\end{choices}

\question
Organisms living in the northern rocky intertidal face competing selective forces. Larger individuals benefit by having a larger internal volume, which helps to conserve body temperature. However, smaller individuals are less likely to be dislodged by large waves from storms. Thus, selection may favor intermediate body size phenotype for any particular species. The mode of selection is

\begin{choices}
\choice natural selection. 
\choice phenotype selection. 
\choice disruptive selection. 
\correctchoice stabilizing selection. 
\choice directional selection. 
\end{choices}

\question
In science, a hypothesis is only useful if it

\begin{choices}
\choice agrees with the conclusion.
\choice is proven correct.
\correctchoice can be tested.
\choice can be proven correct.
\choice explains the natural world.
\end{choices}

\question
The overuse of drugs such as penicillin has caused an increase in the number of bacterial species that are resistant to the drugs. The most likely explanation for this result is:

\begin{choices}
\choice Bacteria began making proteins resistant to the drugs. 
\correctchoice A few bacteria were naturally resistant and natural selection favored their increase. 
\choice Most bacteria can resist drugs. 
\choice The drugs caused mutations in the bacteria that gave them resistance. 
\choice Penicillin was not an effective drug. 
\end{choices}

\question
According to Darwin's concept of descent with modification:

\begin{choices}
\choice All organisms overproduce offspring, leading to differences in survivorship. 
\choice All species have evolved from a single common ancestor. 
\choice All organisms overproduce offspring, leading to different adaptations. 
\choice Individuals with greater fitness will evolve more quickly than individuals with lower fitness. 
\correctchoice Species evolve from ancestral forms through the accumulation of different adaptations. 
\choice Populations evolve by natural selection acting on individuals. 
\end{choices}

\newpage

\question
Which of these conditions should completely prevent the occurrence of natural selection in a population over time?

\begin{choices}
\choice The population size is large. 
\choice The population lives in a habitat where there are no competing species present. 
\choice The environment is changing at a relatively slow rate. 
\correctchoice All variation between individuals is due only to environmental factors. 
\choice The population mates non-randomly.
\end{choices}

\question
Violation of these two Hardy-Weinberg assumptions always causes homozygosity to increase.

\begin{choices}
\choice Natural selection and gene flow.
\choice Genetic drift and natural selection.
\choice Non-random mating and gene flow.
\correctchoice Genetic drift and non-random mating.
\choice Mutation and natural selection.
\choice Gene flow and mutation.
\end{choices}
